\section{Completion as a categorical construction with arithmetic interfaces}
\label{sec:cm-completion}

The existing repository already has a substantial abstract interface.  In
\path{Foundations/SetoidInterpretation.lean}, numerical domain and equality
formulas interpret a source theory, and the theorem \texttt{sound} transports
its intuitionistic derivations once closure, congruence, and interpreted axioms
have been proved.  \texttt{PresentedCategory.lean} and \texttt{Yoneda.lean}
supply actual arithmetic derivations for specified categorical presentations.
Consequently the missing step is usually not the equational manipulation of a
category.  It is constructing a suitable presentation of its continuous
objects and proving that its proposed operations preserve those objects.
This section supplies one such step, without asserting that every analytic
category has a numerical universe.

\subsection{Two levels that must remain distinct}

A \emph{rational metric presentation} consists of a decidable set of numerical
codes $D$, a distinguished valid code, and a fixed primitive recursive rational
pseudometric $d:D\times D\to\Q_{\geq0}$, with proofs of its pseudometric laws.
The use of a pseudometric permits different rational codes to denote the same
point.  Rational arithmetic means coded integer numerators and positive
natural denominators; all inequalities below can be cleared of denominators.
There is no assumption that $D$ is finite or discrete in its metric.

Externally, its completion $\widehat D$ has the usual mathematical meaning.
Internally, a family of names is given by a \emph{specified} PR function
$q(p,n)\in D$, where $p$ is a finite tuple of numerical parameters.  Its
validity is an arithmetic formula in $p$.  A different producer $q'$ is a
different instance of the scheme, unless a fixed uniform PR family containing
both has actually been supplied.  We do not form a numerical sort of all PR
programs equipped with a total PR evaluator.  The repository's
\texttt{UniversalEvaluatorObstruction.lean} is directly relevant to this
restriction.

\begin{definition}[Fast names and equality]
\label{def:cm-fast}
A fast name is a sequence $q_n\in D$ satisfying
\[
             n\leq m\quad\Longrightarrow\quad d(q_n,q_m)\leq 2^{-n}.
\tag{C1}\label{eq:cm-fast}
\]
Two fast names are equivalent if
\[
 q\sim r\quad\Longleftrightarrow\quad
 \forall k\ \exists N\ \forall n\geq N\quad d(q_n,r_n)\leq 2^{-k}.
\tag{C2}\label{eq:cm-eq}
\]
For fixed PR producers these are first-order arithmetic formulas, not
quantification over functions.
\end{definition}

Reflexivity and symmetry of $\sim$ follow from the metric laws.  For
transitivity, instantiate the two equivalences at $k+1$, take the maximum of
the resulting two indices, and apply the triangle inequality.  This uses
intuitionistic existential elimination and introduction and contains no choice
principle.  Constant names embed $D$.  The rational sequence
$d(q_{n+1},r_{n+1})$ is fast, since the reverse triangle inequality gives, for
$m\geq n$,
\[
 |d(q_{n+1},r_{n+1})-d(q_{m+1},r_{m+1})|
 \leq 2^{-n-1}+2^{-n-1}=2^{-n}.
\]
It therefore supplies the completed distance in the semantic interpretation.
A fast name is at distance at most $2^{-n}$ from its $n$th constant
approximation.  These observations explain the completion notation.  They do
not imply that an arbitrary sequence of PR names has a PR limit; a uniform
producer and an effective outer Cauchy rate are needed for that assertion.

\subsection{A dense evaluator extends with one precision substitution}

Let $D_X,D_Y$ be two presentations.  A dense evaluator is a fixed PR function
$E:D_X\times\N\to D_Y$.  Its output need only approximate the value at a dense
input; square roots, special functions, and bounded linear operators can be
handled this way.  Let $\mu:\N\to\N$ be a specified PR function, nondecreasing
and satisfying $\mu(k)\geq k$.  The essential hypothesis is the following
finite, rational inequality:
\[
\begin{split}
 d_X(a,b)\leq 2^{-\mu(k)},\quad r,s\geq k+2
 \quad\Longrightarrow\quad
 d_Y(E(a,r),E(b,s))\leq2^{-k}.
\end{split}
\tag{C3}\label{eq:cm-joint}
\]
The theorem uses a proof of this implication; it does not test its universal
validity.  When $E(a,r)$ is known to approximate a map $f(a)$ within $2^{-r}$
and $f$ has spatial modulus $\mu_0$, the normalized choice
\[
     \mu(k)=\max\{k,\mu_0(0),\ldots,\mu_0(k+1)\}
\]
suffices: the three contributions are $2^{-k-2}$, $2^{-k-1}$,
and $2^{-k-2}$.  Thus \eqref{eq:cm-joint} is a convenient way to combine
spatial continuity and evaluation error.

\begin{theorem}[Dense extension with explicit demand]
\label{thm:cm-extension}
Under \eqref{eq:cm-joint}, define
\[
             (\widehat E q)_n=E(q_{\mu(n)},n+2).
\tag{C4}\label{eq:cm-extension}
\]
This is a fast name.  Equivalent input names give equivalent output names.
For a specified PR input family it is a specified PR output family, and both
assertions are first-order schemes over the two families.  Semantically this
construction is a continuous extension to the completions, and it is the
unique uniformly continuous extension of its values on the dense points.
\end{theorem}

\begin{proof}
If $m\geq n$, monotonicity gives $\mu(m)\geq\mu(n)$, so
\eqref{eq:cm-fast} gives
$d_X(q_{\mu(n)},q_{\mu(m)})\leq2^{-\mu(n)}$.
The output precisions $n+2,m+2$ satisfy the two guards in
\eqref{eq:cm-joint}.  Its instance at $k=n$ proves
\[
      d_Y((\widehat E q)_n,(\widehat E q)_m)\leq2^{-n}.
\]
This is the whole fast-name calculation.

Now suppose $q\sim r$, and fix an output challenge $k$.  Use
\eqref{eq:cm-eq} at input precision $\mu(k)$, obtaining $N$.  For
$n\geq\max(N,k)$, the inequality $\mu(n)\geq n$ makes
$d_X(q_{\mu(n)},r_{\mu(n)})\leq2^{-\mu(k)}$.  Both evaluator precisions
are $n+2\geq k+2$.  Equation \eqref{eq:cm-joint} at the fixed challenge $k$
therefore gives the required output distance $\leq2^{-k}$.  This proves
extensionality without selecting a function of the existential witnesses in
\eqref{eq:cm-eq}.

For continuity in the semantic completion, give the input distance a spare
factor of two.  If
$d_{\widehat X}(x,y)\leq2^{-\mu(k+1)-1}$, sufficiently late approximations
have distance at most $2^{-\mu(k+1)}$.  Applying
\eqref{eq:cm-joint} to those approximations and passing to the named limit
gives output distance at most $2^{-k-1}$, hence at most $2^{-k}$.  Thus
$k\mapsto\mu(k+1)+1$ is a safe modulus for the extension.  This semantic
sentence can equally be replaced by its finite rational inequalities.

For uniqueness, suppose $F,G:\widehat X\to\widehat Y$ are uniformly
continuous and agree on constant names.  Given $x=[q]$ and $k$, choose a dense
approximation $q_N$ close enough to $x$ for both output errors to be at most
$2^{-k-1}$.  Then
\[
 d(Fx,Gx)\leq d(Fx,Fq_N)+d(Fq_N,Gq_N)+d(Gq_N,Gx)\leq2^{-k}.
\]
The middle term is zero.  Since this holds at every precision, the two output
names are equivalent.  For specified moduli, $N$ is the maximum of their
explicit input demands; no compactness argument is hidden here.
\end{proof}

For exact dense maps $f:D_X\to D_Y$ with modulus $\mu_f$, the simpler lift
$(\widehat f q)_n=f(q_{\mu_f(n)})$ is already fast.  The identity has
$\mu_{\mathrm{id}}(n)=n$.  For $g\circ f$, use
$\mu_{g\circ f}=\mu_f\circ\mu_g$; the two composite lifted expressions are
literally identical.  Approximate dense evaluators compose by substitution of
the corresponding name transformations, and uniqueness identifies different
valid schedules for the same dense map.  In this precise external or
schematic sense completion is functorial.  It is not an internal category
whose objects are all possible PR implementations.

\begin{proposition}[Products and change of dense presentation]
\label{prop:cm-products-density}
The completion construction preserves finite products equipped with the
maximum metric.  Moreover, suppose $f:D_X\to D_Y$ is a specified PR isometry
and a PR selector $s:D_Y\times\N\to D_X$ satisfies
\[
                    d_Y(f(s(b,k)),b)\leq2^{-k}.
\tag{C5}\label{eq:cm-dense-selector}
\]
Then $\widehat f$ is an isomorphism of the presented completions, with an
explicit inverse on names.
\end{proposition}

\begin{proof}
A name in the maximum product metric is fast exactly when both coordinate
names are fast.  Equivalence in the product is equivalent to equivalence in
both coordinates: the forward implication follows by projection; for the
reverse implication take the maximum of the two eventual indices.  Pairing
and projection therefore give inverse name transformations.  The terminal
one-point presentation covers the empty product.

For the second assertion, define
\[
                  (Hy)_n=s(y_{n+2},n+2).
\tag{C6}\label{eq:cm-density-inverse}
\]
For $m\geq n$, the isometry law and three applications of the triangle
inequality give
\[
\begin{split}
 d_X((Hy)_n,(Hy)_m)
 &=d_Y(f(Hy)_n,f(Hy)_m)\\
 &\leq2^{-n-2}+d_Y(y_{n+2},y_{m+2})+2^{-m-2}\\
 &\leq3\,2^{-n-2}\leq2^{-n}.
\end{split}
\]
Here $f(Hy)_n$ denotes $f((Hy)_n)$.  Thus $H$ produces a fast name.  Also
$d_Y(f((Hy)_n),y_{n+2})\leq2^{-n-2}$, so $\widehat f(Hy)\sim y$.
For an $X$-name $x$, the isometry law similarly gives
$d_X((H\widehat f x)_n,x_{n+2})\leq2^{-n-2}$, hence
$H\widehat f x\sim x$.  Finally, if $y\sim z$, applying the triangle
inequality to $f((Hy)_n),y_{n+2},z_{n+2},f((Hz)_n)$ proves $Hy\sim Hz$.
All output indices in these estimates are displayed; for example an
$n$-stage error bounded by $5\,2^{-n-2}$ is below $2^{-k}$ once
$n\geq k+1$.  Thus the two inverse equations are first-order equality
statements about the supplied name families.
\end{proof}

Condition \eqref{eq:cm-dense-selector} is positive information.  An unadorned
assertion that $f(D_X)$ is dense need not provide a PR approximation selector.
This is the continuous analogue of the finite-cover versus cardinality-bound
distinction already identified in the repository's categorical packet.
The proposition permits safe transport between, for example, rational and
dyadic descriptions of a compact interval, or two explicit dense descriptions
of a normed space, once the displayed selector is proved.

\subsection{A direct Lean implementation boundary}

The first implementation should keep codes, rational distances, and producers
explicit.  One can use the following metatheoretic records and then replace
the fields asserting inequalities by the corresponding derivations in
\texttt{Prf PRAAx}:
\begin{verbatim}
structure MetricPresentation where
  valid : Nat -> Bool
  dist  : Nat -> Nat -> RatCode
  -- Proofs: closure, nonnegativity, symmetry, triangle.

def Fast (D : MetricPresentation) (q : Nat -> Nat) : Prop :=
  (forall n, D.valid (q n) = true) /\
  (forall n m, n <= m -> D.dist (q n) (q m) <= dyadic n)

structure DenseEvaluator (X Y : MetricPresentation) where
  eval : Nat -> Nat -> Nat
  mu   : Nat -> Nat
  -- PR realizers and monotonicity proofs are separate fields.
  joint : forall a b k r s, ValidInputs a b ->
    X.dist a b <= dyadic (mu k) -> k + 2 <= r -> k + 2 <= s ->
    Y.dist (eval a r) (eval b s) <= dyadic k

def extend (E : DenseEvaluator X Y) (q : Nat -> Nat) (n : Nat) :=
  E.eval (q (E.mu n)) (n + 2)
\end{verbatim}

This block is a blueprint, not compilable source and not an assertion that
\texttt{Nat -> Nat} is an object-language sort.  In the arithmetic layer it is
replaced by an explicit PR term $q(p,n)$ and by formulas
$\mathrm{Fast}_q(p)$ and $\mathrm{Eq}_{q,r}(p,p')$.  The principal target is
\[
 \ipra\vdash\forall p\bigl(\mathrm{Fast}_q(p)
                  \to\mathrm{Fast}_{\widehat E q}(p)\bigr),
\]
with the modulus and joint-law hypotheses supplied as derivations.  Its proof
instantiates those hypotheses at arbitrary terms and combines rational
inequalities; it does not induct on $\mathrm{Fast}_q$, which contains
unbounded quantifiers.  Product and density-selector targets use the same
procedure.  The existing setoid-interpretation theorem then applies to a
\emph{fixed} numerical family closed under the operations that have actually
been defined.

\section{A finite-observation semantics for compact metric structures}
\label{sec:cm-model}

Continuous model theory uses metric structures, uniformly continuous symbols,
real-valued predicates, and $\inf/\sup$ quantifiers; this is the setting
surveyed by Ben Yaacov, Berenstein, Henson, and Usvyatsov
\cite{bbhu2008}.  The following development supplies a deliberately narrower
arithmetic realization with explicit total-boundedness data.  It preserves
quantifier order, including alternation.  It neither imports the general
compactness theorem of continuous logic nor replaces it by a finite-model
claim.

\subsection{Language, moduli, and finite restrictions}

Fix finitely many metric sorts and a finite signature. In the exact dense
version, constants have supplied dense codes; named constants are covered
by the approximate-evaluation variant below. For simplicity first suppose
that function symbols act by fixed
PR maps between the dense presentations and that the predicates are fixed
PR rational functions taking values in $[0,1]$.  Each symbol is accompanied
by a PR modulus and proofs of its continuity inequalities.  The dense
extension theorem interprets the symbols on the completions.

Terms are variables, constants, and finite applications of function symbols.
An atomic expression is a predicate applied to terms, including the bounded
metric predicate.  Formulas are built from atoms and rational constants in
$[0,1]$ by $\min$, $\max$, $1-u$, and the quantifiers
$\inf_{x:s}$ and $\sup_{x:s}$.  With the convention that $0$ is true,
$\max$ plays the role of conjunction and $\min$ that of disjunction.
The operation $1-u$ is an arithmetic connective in the object being studied;
its use does not add classical excluded middle to the surrounding iPRA
proof.

A modulus of a term or atom is computed by composition.  For a variable use
$\alpha(k)=k$.  If $f$ has modulus $\mu_f$ and the arguments have moduli
$\alpha_i$, use
\[
            \alpha_{f(t_1,\ldots,t_r)}(k)
                    =\max_i\alpha_i(\mu_f(k)).
\tag{M1}\label{eq:cm-atom-modulus}
\]
Use the same rule for a predicate.  The metric has safe modulus $k+1$ with
respect to the maximum distance on its two arguments.  Constants have no
input demand.  Extend the modulus trivially to unused variables.  For a fixed
formula $\varphi$, let $\alpha_\varphi(k)$ be the maximum of the moduli of
all its atomic occurrences, normalized to be nondecreasing and at least $k$.
Then $\delta=2^{-\alpha_\varphi(k)}$-close assignments make every atomic
value differ by at most $2^{-k}$.

For each sort choose a finite nonempty list $A_s$ of dense points.  The finite
restriction $V_A(\varphi,\bar a)\in\Q\cap[0,1]$ is defined recursively:
use the dense evaluations at atoms, the indicated rational operations at
connectives, and the finite minimum or maximum over $A_s$ at a quantifier of
sort $s$.  Bound variables have their usual scope; the recursion does not
interchange two quantifiers.  Terms continue to be evaluated in the ambient
dense presentation.  Thus $A_s$ need not be closed under a function symbol.
In particular $V_A$ is not being advertised as a finite submodel of the
original structure.

\begin{theorem}[Quantitative back-and-forth for finite restrictions]
\label{thm:cm-backforth}
Let $A_s,B_s$ be finite nonempty lists for every sort.  Suppose each point of
$A_s$ is within $\delta$ of some point of $B_s$, and conversely.  Suppose
also that corresponding free-variable assignments are coordinatewise within
$\delta$, and that every atomic occurrence of $\varphi$ changes by at most
$\eta$ on $\delta$-close assignments.  Then
\[
                |V_A(\varphi,\bar a)-V_B(\varphi,\bar b)|\leq\eta.
\tag{M2}\label{eq:cm-transfer}
\]
The error does not acquire a factor from the number of quantifier alternations.
If atomic values are themselves evaluated with errors at most $e_A,e_B$, the
right side becomes $\eta+e_A+e_B$.
\end{theorem}

\begin{proof}
We use structural induction on the fixed formula to produce its arithmetic
proof.  The atom case is the hypothesis.  For $1-u$, the error is unchanged.
For $\min$ and $\max$, the elementary inequality
\[
 |\max(u_1,u_2)-\max(v_1,v_2)|
 \leq\max\{|u_1-v_1|,|u_2-v_2|\}
\]
and its minimum counterpart preserve the same bound.  Each follows by the
two inequalities $u_i\leq v_i+\eta$ and $v_i\leq u_i+\eta$; taking the
appropriate maximum or minimum proves both directions.

Consider $\inf_{x:s}\psi$.  For any $a\in A_s$, choose a $b\in B_s$
within $\delta$.  The extended assignments still satisfy the same
coordinatewise relation.  The induction hypothesis gives
\[
 V_A(\psi,\bar a,a)\geq V_B(\psi,\bar b,b)-\eta
             \geq\min_{b'\in B_s}V_B(\psi,\bar b,b')-\eta.
\]
Taking the finite minimum over $a$ proves one inequality.  The reverse
covering direction proves the other, with $A$ and $B$ exchanged.  Alternatively,
a first minimizing entry can be selected by a bounded scan and used as the
witness.  For $\sup$, replace minima by maxima and reverse the inequalities.
The error at a child is not added to a new rounding error: the corresponding
bound-variable choices were already included in its atom-level hypothesis.
This explains the absence of a depth factor.

Finally, evaluating every atom within an error $e$ changes the entire finite
formula by at most $e$, by the identical induction with the same domains.
Applying this on both sides of \eqref{eq:cm-transfer} gives
$e_A+\eta+e_B$.  Only rational operations and finite minima, maxima, and
searches appear in this proof.
\end{proof}

The finite covering relation is a bounded formula: after fixing the lists,
all its existential witnesses have an explicit list-length bound.  It may be
Booleanized.  Induction used to compute a minimum is on a finite prefix with
its rational lower-bound invariant.  Thus the proof does not require induction
on a formula involving quantification over an entire completion.  An external
induction on the fixed source syntax generates a finite iPRA proof, much as
in the existing interpretation theorem.

\subsection{From nets to names of formula values}

Assume each sort has a fixed PR net producer $N_s(t)$, of positive PR length,
with the supplied covering theorem
\[
       \forall a\in D_s\ \exists j<|N_s(t)|\quad
                   d_s(a,N_s(t)_j)\leq2^{-t}.
\tag{M3}\label{eq:cm-nets}
\]
The unbounded quantifier over $a$ in this theorem is a hypothesis already
proved; validity of a purported net producer is not decided by inspecting one
finite list.  This is a quantitative total-boundedness assumption, stronger
than merely calling a sort bounded.

\begin{theorem}[PR values for compactly presented metric formulas]
\label{thm:cm-values}
Fix the signature, its PR producers, and a formula $\varphi$.  For specified
fast input names $q^1,\ldots,q^r$, put
\[
 v_n=V_{N(\alpha_\varphi(n))}
       \bigl(\varphi,q^1_{\alpha_\varphi(n)},\ldots,
                          q^r_{\alpha_\varphi(n)}\bigr).
\tag{M4}\label{eq:cm-formula-name}
\]
Then $v_n$ is a PR fast Cauchy name.  Its represented value is independent of
valid input names and of the chosen net producers.  In the semantic
completion it gives the usual continuous $\inf/\sup$ interpretation of the
formula.  Each instance has a first-order iPRA realization by finite
inequalities.
\end{theorem}

\begin{proof}
For $m\geq n$, the free-variable tuples in \eqref{eq:cm-formula-name} are
within $\delta=2^{-\alpha_\varphi(n)}$ by fastness and monotonicity.  Every
point of the coarser net is a dense point, so the finer net covers it within
$2^{-\alpha_\varphi(m)}\leq\delta$; the coarser net covers every point of
the finer net within $\delta$ as well.  Theorem~\ref{thm:cm-backforth} gives
$|v_n-v_m|\leq2^{-n}$.  This is an explicit Cauchy proof; an independently
assumed existence of a real infimum is unnecessary.

The same comparison proves independence from two choices of net producers.
For equivalent input names, fix an output precision, instantiate the
input equivalences at the required atomic precision, and take stages beyond
the maximum of the resulting finitely many indices.  Applying
Theorem~\ref{thm:cm-backforth} at that fixed precision shows that the two
output names are equivalent.

To identify the quantifier semantics, it is enough to handle an infimum.
Fix the other variables.  For any dense $a$, the net theorem supplies points
within the required atomic modulus of $a$; their formula values bound the
finite minimum from above up to the displayed error.  Hence the named value
is a lower bound for all dense-input values.  Conversely every finite
minimum is attained at a net entry.  Those entries give values arbitrarily
close above the named lower bound.  Continuity extends the lower-bound
inequality from the dense points to arbitrary named points.  This is exactly
the defining lower-bound and approximate-attainment property of the infimum.
The supremum argument is dual.  Notice that the approximate minimizing net
entry may change with precision.  No coherent convergent sequence of
minimizers has been asserted.

Every evaluator in \eqref{eq:cm-formula-name} uses fixed PR symbols and a
finite syntax tree.  Evaluation at a quantifier is a bounded scan through a
PR-length list; nested scans have a fixed finite nesting depth.  Consequently
the displayed producer is PR.  The arithmetic proof applies the covering
theorems only to the particular finite entries used by the scans, and uses
bounded invariants for those scans.  It does not induce on a real-valued
satisfaction predicate.
\end{proof}

A useful version permits PR rational approximations to dense atomic values.
Evaluate all atoms to error $2^{-n-2}$, use nets and free-variable stages at
$\alpha_\varphi(n+2)$, and retain the same finite evaluation.  Comparing
stages $n\leq m$ gives an error at most
\[
                2^{-n-2}+2^{-n-2}+2^{-m-2}\leq 2^{-n}.
\]
Thus the conclusion includes predicates defined by analytic evaluation
algorithms, not only rational polynomials.  For an approximate dense function
symbol, first compose its evaluator with the finitely many surrounding term
and predicate evaluators, allocating precision along that fixed term tree.
The resulting atomic evaluator and atomic modulus are the explicit inputs to
the theorem.

The theorem gives a practical model-theoretic transfer principle.  If two
compact presentations have uniformly close finite observations, and
corresponding atomic evaluations satisfy the required error bounds, then
every fixed continuous formula has the same controlled error.  If all errors
tend to zero with a supplied PR rate, their sentence values agree.  This is a
restricted, quantitatively witnessed form of elementary equivalence.  It
contains no assertion that every finitely satisfiable theory has a model.

\subsection{A nontrivial alternating-quantifier example}
\label{subsec:cm-minimax}

On the compact sort $[0,1]$, consider the rationally evaluated atom
\[
                  P(c,x)=|x^2-cx|
\]
and the continuous sentence
\[
             v=\inf_{c\in[0,1]}\ \sup_{x\in[0,1]} P(c,x).
\tag{M5}\label{eq:cm-example}
\]
Its quantifier order is essential: $c$ is chosen once before the adversarial
choice of $x$.  The atom takes values in $[0,1]$, and if both coordinates
change by at most $\delta$, then
\[
 |P(c,x)-P(c',x')|
 \leq |x^2-x'^2|+|cx-c'x'|\leq4\delta.
\]
Hence dyadic nets of mesh $2^{-k-2}$ give a formula-value error at most
$2^{-k}$.  Exact finite rational minimax calculations therefore define the
value in iPRA, without a real optimization oracle.

The semantic value and optimizer can also be identified explicitly.  For
$0\leq c\leq1$,
\[
       x^2-cx=(x-c/2)^2-c^2/4,
\]
so its negative part has maximum $c^2/4$, attained at $x=c/2$.
Its positive part has maximum $1-c$, attained at $x=1$: indeed
\[
       (1-c)-(x^2-cx)=(1-x)(1+x-c)\geq0.
\]
Consequently
\[
        \sup_x P(c,x)=\max\{c^2/4,1-c\}.
\]
The first expression increases and the second decreases.  Their common value
occurs at $c_*=2\sqrt2-2\in[0,1]$, and is
\[
                         v=3-2\sqrt2.
\tag{M6}\label{eq:cm-minimax-value}
\]
The following order cases may be performed first on rational $c$: its side
of the algebraic number $c_*$ is determined by the rational sign test on
$(c+2)^2-8$.  The factorization
$(c+2)^2-8=(c-c_*)(c+c_*+4)$ has a positive second factor.
The resulting non-strict estimates then extend to arbitrary named $c$ by
continuity.  Thus no exact sign test on an arbitrary real name is required.
For $c\leq c_*$ the second term is at least $1-c_*$; for $c\geq c_*$ the
first term is at least $c_*^2/4$.  This proves global minimality, and the strict
versions prove uniqueness.  The occurrence of $\sqrt2$ can be interpreted by
a fixed rational square-root name with its elementary inequalities.  All the
displayed algebraic estimates are denominator-cleared arithmetic statements.
The accompanying reference code computes the finite restricted sentence and
its minimizing net entry with exact integers and rational output.

In this example a quantitative uniqueness estimate upgrades the finite
witnesses to a coherent optimizer name.  Write
$w(c)=\max\{c^2/4,1-c\}$.  Since $c_*\geq1/2$, the two cases give
\[
                 w(c)-v\geq\tfrac14|c-c_*|.
\tag{M7}\label{eq:cm-minimax-gap}
\]
For $c\leq c_*$ use $1-c=v+c_*-c$; for $c\geq c_*$ use
$(c^2-c_*^2)/4=(c+c_*)(c-c_*)/4\geq(c-c_*)/4$.
Let $c_j$ be the first minimizing entry of the dyadic restriction with mesh
$\delta=2^{-j}$.  The generic formula bound gives $|v_j-v|\leq4\delta$,
and the inner supremum at $c_j$ differs from its finite maximum by at most
$4\delta$.  Hence $w(c_j)\leq v+8\delta$, so
$|c_j-c_*|\leq32\,2^{-j}$.  Therefore $a_n=c_{n+6}$ satisfies
$|a_n-c_*|\leq2^{-n-1}$, and
$|a_n-a_m|\leq2^{-n}$ for $m\geq n$.
The first-minimum scan thus constructs a PR fast name of the exact optimizer.
The additional ingredient is the proved gap \eqref{eq:cm-minimax-gap}, not
an assumption that arbitrary approximate minimizers are coherent.

\subsection{Sharp boundaries and the remaining Lean work}

There are three limitations with mathematical content.
First, boundedness is not enough for \eqref{eq:cm-nets}.  For example, the unit
ball of an infinite-dimensional Hilbert space has infinitely many orthogonal
unit vectors separated by a fixed positive distance.  No finite sufficiently
small net covers it.  The compact formula evaluator therefore cannot be
applied to arbitrary quantification over that unit ball.  A specified compact
subset with a uniform tail bound, a finite-rank restriction, or a different
witness contract is needed.

Second, a name of a formula value is not a decision procedure for whether it
is zero.  For a uniform example, let $h(e,t)$ be the fixed PR test that a
machine $e$ first halts at step $t\geq1$, with a harmless one-step delay
ensuring no halt at step zero.  On a one-point metric sort, give a predicate
the value $r_e=2^{-t}$ if such a first halting time exists, and $0$ otherwise.
Its $k$th rational approximation is $2^{-t}$ if a first halt is found with
$t\leq k$, and $0$ otherwise.  Bounded simulation is PR and the approximation
error is at most $2^{-k-1}$.  Its spatial modulus is trivial, and its sort has
a one-point net at every precision.  A total algorithm deciding whether
$r_e=0$ would decide nonhalting and hence halting.  Thus even this simplest
compact metric semantics cannot supply uniform exact truth tests.  This is a
metatheoretic obstruction to a proposed uniform procedure, not an
independence claim about any particular named sentence.

Third, the finite minima used above supply approximate witnesses.  They do not
supply a coherent exact witness unless an additional mechanism, such as a
contraction or a quantitative uniqueness modulus, is available.  Likewise
this construction provides no ultrafilter, saturated model, or model of an
arbitrary consistent theory.

For Lean, the principal new definitions are a fixed finite multisorted
signature, its term and continuous-formula syntax, finite environments,
rational finite evaluation, the two-sided finite covering relation, and
per-atom modulus contracts.  The first central lemma is
Theorem~\ref{thm:cm-backforth}; its proof follows the displayed syntax cases,
with the quantifier case reduced to a lemma about minima of two finite lists
related in both directions.  The second central lemma is
Theorem~\ref{thm:cm-values}, obtained by substituting two net stages and two
name stages into that result.  Approximate attainment is the first minimizing
entry from the finite scan.  No real-number library is needed for these
arithmetic endpoints; a separate semantic interpretation can identify the
result with conventional continuous logic.

The delivered \texttt{QuantitativeKernel.lean} checks a qualitative nucleus:
positive formulas built from atoms, conjunction, disjunction, and the two
quantifiers transfer across a relation with selected forward and backward
lifts, provided atoms transfer.  Its \texttt{positive\_transfer} theorem
contains no quantifier-depth loss; \texttt{eval\_correct} connects finite
Boolean evaluation to the corresponding finite-domain semantics.  These are
Lean metatheorems.  They are not a completed formalization of the numerical
continuous-logic estimates above, nor are they silently identified with
\texttt{Prf PRAAx} derivations.  To reach that endpoint, instantiate the finite
minimum lemmas in the repository's rational arithmetic, Booleanize their
bounded list invariants, and construct the syntax-indexed family of
\texttt{Prf PRAAx} proof terms.  The existing
\texttt{SetoidInterpretation.sound} provides the subsequent interpretation
interface for numerical sort families whose closure has been established.
